\subsection{General remarks}

Let A be a subset of a topological space E. For a sequence \((x_{n})_{n\in\mathbf{N}}\) of points of A to have a point x of E as a limit point, it is necessary and sufficient that the following condition is satisfied (GT, I, § 7, No.~3):

\begin{enumerate}
\def\labelenumi{(\Alph{enumi})}
\tightlist
\item
  For every integer \(m \geqslant 0\) and every neighbourhood \(\mathbf{U}\) of \(x\) , there exists an integer \(n \geqslant m\) such that \(x_{n} \in \mathbf{U}\) .
\end{enumerate}

A sequence of the form \((y_{k})_{k\in\mathbb{N}}\) with \(y_{k}=x_{n_{k}}\) for a strictly increasing sequence \((n_{k})_{k\in\mathbb{N}}\) of positive integers is called an extracted sequence of the sequence \((x_{n})_{n\in\mathbb{N}}\) . If there exists an extracted sequence of the sequence \((x_{n})_{n\in\mathbb{N}}\) which converges to x, then x is a limit point of \((x_{n})\) ; conversely, if x has a countable fundamental system of neighbourhoods, and x is the limit point of the sequence \((x_{n})\) , then there exists an extracted sequence of \((x_{n})\) converging to x.

On account of GT, IX, § 2, No.~9, corollary, we conclude that when E is metrizable, the following conditions are equivalent :

\begin{enumerate}
\def\labelenumi{\alph{enumi})}
\item
  the set A is relatively compact in E;
\item
  every infinite sequence of points of A has a limit point in E;
\item
  from every infinite sequence of points of A, we can extract a sequence which converges to a point of E.
\end{enumerate}

In this section, we shall extend this criterion to certain non metrizable topological vector spaces. The following proposition enables us to reduce the study of compact sets to that of weakly compact sets in a number of cases.

PROPOSITION 1. --- Let E be a Hausdorff locally convex space, and A a subset of E. Let \(E_{\sigma}\) denote the space E with the weakened topology.

\begin{enumerate}
\def\labelenumi{\alph{enumi})}
\item
  If every infinite sequence of points of \(\mathbf{A}\) has a limit point in \(\mathbf{E}\) , then \(\mathbf{A}\) is precompact in \(\mathbf{E}\) .
\item
  In order that A be relatively compact in E, it is necessary and sufficient that it is precompact in E and relatively compact in \(E_{\sigma}\) .
\end{enumerate}

We shall prove \(a)\) by reductio ad absurdum. If \(A\) is not precompact, then by th. 3 of GT, II, § 3, No.~7, it follows that there exists a symmetric convex neighbourhood \(V\) of 0 in \(E\) such that \(A\) cannot be covered by a finite number of translates of \(V\) .

In other words, if \(x_{0}, x_{1}, \ldots, x_{n-1}\) are points of A, then \(A \not\subset \bigcup_{0 \leqslant i < n} (x_{i} + V)\) and so there exists a point \(x_{n}\) of A such that \(x_{n} - x_{i} \notin V\) for \(0 \leqslant i < n\) . Then, by induction on the integer n, we can construct an infinite sequence \((x_{n})_{n \in \mathbb{N}}\) of points of A such that \(x_{n} - x_{m} \notin V\) whenever n \textgreater{} m; since V is symmetric, we also have \(x_{m} - x_{n} \notin V\) for \(m \neq n\) and the sets \(x_{n} + \frac{1}{2}V\) are disjoint. For every point x in E, there exists at most one integer \(n \geqslant 0\) such that \(x_{n} \in x + \frac{1}{2}V\) , hence the sequence \((x_{n})_{n \in \mathbb{N}}\) does not have any limit point. This proves a).

Now suppose that A is precompact in E and is contained in a compact subset B of \(E_{\sigma}\) . Then B is complete in \(E_{\sigma}\) , hence also in E (IV, p.~5, Remark 2). We have \(\overline{A} \subset B\) , hence A is relatively compact in E. The converse is evident and b) follows.

